\subsection{Problem statement}
\subsubsection{Multivariable system}
Consider the multivariable system
\begin{equation}
    \dot{x} = f(t, x) + G(t, x)u
    \label{eq:moreno-system}
\end{equation}
where $x \in \R^n$ is the state vector, $u \in \R^n$ is the contorl input, $f: \R_+ \times \R^n \rightarrow \R^n$ is an uncertain vector field containing uncertainties and/or perturbations, and $G: \R_+ \times \R^n \rightarrow \R^{n \times n}$ is the uncertain input matrix.

\subsubsection{Input matrix}
The uncertain input matrix can be rewritten
\begin{equation}
    G(t, x) = \left(I + \Delta_G(t, x)\right)G_0(t, x)
\end{equation}
where $G_0: \R_+ \times \R^n \rightarrow \R^{n \times n}$ corresponds to the nominal part, which is assumed to be invertible, and $\Delta_G: \R_+ \times \R^n \rightarrow \R^{n \times n}$ is the uncertain term

\subsection{Multivariable super-twisting algorithm}
\subsubsection{Control law by Moreno et al. \cite{Mor:21}}
The MGSTA is given by
\begin{equation}
    \begin{split}
        u &= -k_1\phi_1(x) + bG_0^{-1}(t, x)v \\
        \dot{v} &= -k_2\phi_2(x)
    \end{split}
    \label{eq:moreno-controller}
\end{equation}
The nonlinear functions $\phi_i: \R^n \rightarrow \R^n$ are monotonically increasing, $\phi_1$ is continuous everywhere and $\phi_2$ is continuous everywhere except possibly at $x=0$.


Nonlinear functions
\begin{align}
    \phi_1 &= \left(\alpha\norm{x}^{-p}+\beta\right)x \label{eq:moreno-phi_1}\\
    \phi_2 &= J(x)\phi_1(x) = c(x)\phi_1(x) \label{eq:moreno-phi_2}
\end{align}



\subsubsection{Decomposed perturbations}
The uncertain term $f(t,x)$ is assumed to be decomposed as
\begin{equation}
    f(t, x) = f_1(t,x) + f_2(t,x) = \Delta_1(t,x)\phi_1(x) + f_2(t,x)
    \label{eq:moreno-eq-10}
\end{equation}
where the function $f_2(t,x)$ is such that
\begin{equation}
    \frac{d}{dt}\left[\left(b\left(I + \Delta_G\right)^{-1}f_2\right)\right] = \Delta_2(t,x)\phi_2 + \Delta_3(t,x)\dot{x}
    \label{eq:moreno-eq-11}
\end{equation}
where $\Delta_i : \R_+ \times \R^n \rightarrow \R^{n \times n}, i=1,2,3$.

For the dynamical system \eqref{eq:moreno-system} and controller \eqref{eq:moreno-controller}, \eqref{eq:moreno-phi_1} and \eqref{eq:moreno-phi_2}, the following conditions hold for all $t \geq 0, x \in \R^n$.
\begin{assumption}
    There exit constants $g_M > 0, \delta_g \geq 0$ and $g_m^\Delta > 0$ such that
            \begin{equation}
                \begin{split}
                    \norm{G(t, x)} &\leq g_M \\
                    \norm{\Delta_G(t, x)} &\leq \delta_g \\
                    \text{Sym}\left\{G(t, x)\right\} &\geq g_m^\Delta I
                \end{split}
            \end{equation}
\end{assumption}
\begin{assumption}
    Matrices $\Delta_i, i=1,2,3$, in \eqref{eq:moreno-eq-10} and \eqref{eq:moreno-eq-11}, are bounded, there exist $\delta_i \geq 0$ such that 
            \begin{equation*}
                \norm{\Delta_i(t,x)} \leq \delta_i, i = 1,2,3.
            \end{equation*}
\end{assumption}
\begin{assumption}
    There exist positive constants $0 < \gamma_1^\Delta \leq \gamma_2^\Delta$ such that 
            \begin{equation}
                0 < \gamma_1^\Delta I \leq \text{Sym} \left\{\mathcal{J}(x)\left(I + \Delta_G\right)\right\} \leq \gamma_2^\Delta I
            \end{equation}
            As a consequence of these assumptions, there also exist constants $\gamma_3,\gamma_4,\gamma_5 > 0$, and $0 < \gamma_1^\Delta \leq\gamma_2^\Delta$ such that 
            \begin{equation}
                \begin{split}
                    \text{Gram}\left\{\mathcal{J}G\right\} &\leq \gamma_3 I \\
                    2\text{Sym}\left\{\Delta_G\mathcal{J}G\right\} &\leq \gamma_4 I \\
                    \text{Gram}\left\{\Delta_G^\top\right\} &\leq \gamma_5 I \\
                \end{split}
            \end{equation}
\end{assumption}

\textbf{Closed-loop dynamics}
The closed-loop dynamics including MGSTA is given by
\begin{align}
    \dot{x} &= f_1 + \left(I + \Delta_G\right)\left(-G_0k_1\phi_1 + bz\right) \\
    \dot{z} &= -k_2\phi_2 + \Delta_2\phi_2+\Delta_3 \dot{x}
\end{align}


\begin{theorem}[Moreno et al. \cite{Mor:21}]
Select $\alpha > 0, \beta > 0, b > 0$ and $p \in \left(0, \frac{1}{2}\right]$ and let Assumption 1 be satisfied. Assume further that the inequality 
\begin{equation}
    \gamma_1^\Delta g_m^\Delta > \gamma_4 + 2\sqrt{\gamma_3 \gamma_5}
\end{equation} 
is satisfied. Then, for arbitrary values of $\delta_2, \delta_3, \delta_3$, there exist positive constant controller gains $k_1, k_2$ such that $\left[\sigma^\top, z^\top\right]^\top$ is globally robust finite-stable.
\end{theorem}

\subsubsection{Algorithm 1}
For proper selection of the control gains $k_1$ and $k_2$ consult the following algorithm
\begin{enumerate}
    \item Select $\alpha > 0, \beta > 0, b > 0$ and $p \in \left(0, \frac{1}{2}\right]$. The parameter $p$ should be selecting on the size of the perturbations.
    \item Select $p_2 = p_2^*$ and $p_1 = p_1^*$ such that 
    \begin{equation}
        \gamma_1^\Delta > p_2\mu_4
    \end{equation}
    and 
    \begin{equation}
        \gamma_1^\Delta g_m^\Delta - \gamma_4 - \Gamma_0 (p_1, p_2) > 0
    \end{equation}
    are satisfied for $p_2 = p_2^*$ and $p_1 \geq p_1^*$.  For proper selection of the control gains $k_1$ and $k_2$ consult the following algorithm
    \item Find a value $p^\#_1$ of $p_1$ such that $p_1 \geq p^\#_1$ and
    \begin{equation}
        \Gamma_1(p_1, p_2)p_1 > \Gamma_2(p_1, p_2)
    \end{equation}
    is satisfied for every $p_1 \geq p^\#_1$. This is always feasible.
    For proper selection of the control gains $k_1$ and $k_2$ consult the following algorithm
        \begin{enumerate}
            \setcounter{enumi}{3}
            \item Fix $p_2 = p_2^*, p_1\geq p_1^\#$ and $p_1 p_2 > 1$, choose
            \begin{equation}
                k_2 = b\frac{p_1}{p_2},
            \end{equation}
            and 
            \begin{equation}
                \frac{-\alpha_1 - \bar{\alpha}}{2\alpha_2} < k_1 < \frac{-\alpha_1 + \bar{\alpha}}{2\alpha_2}
            \end{equation}
            with $\bar{\alpha} = \sqrt{\alpha^2_1 - 4\alpha_2 \alpha_0}$
        \end{enumerate}
\end{enumerate}



\subsection{Example: Four-wheeled omidirectional robot}
{Sliding variables, \cite{Mor:21}}
    % \begin{figure}[h]
    %     \includegraphics[width=0.6\textwidth]{figures/moreno-fig1.png}
    % \end{figure}

{Position and orientation, \cite{Mor:21}}
        % \begin{figure}[h]
        %     \includegraphics[width=0.6\textwidth]{figures/moreno-fig2.png}
        % \end{figure}


{Armature voltages, \cite{Mor:21}}
        % \begin{figure}[h]
        %     \includegraphics[width=0.6\textwidth]{figures/moreno-fig3.png}
        % \end{figure}

\subsection{Key takeaways}
{Features}
        \begin{itemize}
            \item Deals well with matched uncertainties
            \item Works for uncertain input gain
            \item Coupled MIMO formulation
        \end{itemize}
{Future work}
        \begin{itemize}
            \item Implement adaptive law
            \item Cleaner formulations and communication
        \end{itemize}

